% ===============================
% Worksheet / Manual Template
% ===============================
\documentclass[12pt, a4paper]{article}

% ===============================
% Metadata
% ===============================
\newcommand*{\CourseCode}{MAT321}
\newcommand*{\CourseTitle}{Real Analysis II}
\newcommand*{\Chapter}{Total Boundness and Compactness}
\newcommand*{\Topics}{$\epsilon$-net and Total Boundness}
\newcommand*{\DocDate}{April 28, 2026}

\include{header}

% ==========================================
% Document
% ==========================================
\begin{document}
\FrontPage
\Large

\begin{heading}
    $\epsilon$-net and Total Boundness
\end{heading}

\vspace{10pt}

\begin{definition}[$\epsilon$-net]
    Let $(X, d)$ be a metric space and $\epsilon > 0$. A subset $N \subseteq X$ is called an \textit{$\epsilon$-net} for $X$ if for every point $x \in X$, there exists a point $n \in N$ such that $d(x, n) < \epsilon$. In other words, the collection of open balls $\{B(n, \epsilon) : n \in N\}$ covers the entire space $X$.
\end{definition}

\clearpage

\begin{problem}[Example of $\epsilon$-net]
    Show that the set $Z = \{-3, -2, -1, 0, 1, 2, 3\}$ is a $1$-net for $\mathbb{R}$.
\end{problem}

\vspace{200pt}

\begin{problem}[Example of $\epsilon$-net]
    The set $B=\{B(m, n) : m, n \in \mathbb{Z}\}$ is a $\sqrt{2}$-net for $\mathbb{R}^2$.
\end{problem}

\clearpage

\begin{definition}[Total Boundness]
    A metric space $(X, d)$ is said to be \textit{totally bounded} if for every $\epsilon > 0$, there exists a finite $\epsilon$-net for $X$. In other words, for any given $\epsilon > 0$, we can find a finite set of points in $X$ such that the open balls centered at those points with radius $\epsilon$ cover the entire space $X$.
\end{definition}

\vspace{300pt} 

\begin{problem}[Example of Total Boundness]
    The interval $[0, 1]$ is totally bounded.
\end{problem}

\clearpage

\begin{problem}[Example of Total Boundness]
    The set $\mathbb{R}$ is not totally bounded.
\end{problem}

\vspace{140pt} 

\begin{problem}[Example of Total Boundness]
    The infinite discrete space $\mathbb{N}$ with the discrete metric is not totally bounded.
\end{problem}

\vspace{140pt} 

\begin{problem}[Example of Total Boundness]
    The infinite-dimensional space $\ell^2$ is not totally bounded.
\end{problem}

\clearpage

\begin{heading}
    Separability and Compactness
\end{heading}

\vspace{10pt}

\begin{theorem}[Totally Boundness and Separability]
    Every totally bounded metric space is separable.
\end{theorem}

\clearpage

\begin{theorem}[Compactness and Separability]
    Every compact metric space is separable.
\end{theorem}

\clearpage

\begin{heading}
    Total Boundness and Compactness
\end{heading}

\vspace{10pt}

\begin{theorem}[Total Boundness and Cauchy Subsequences]
    A metric space $(X, d)$ is totally bounded if and only if every sequence in $X$ contains a Cauchy subsequence.
\end{theorem}

\clearpage

\begin{theorem}[Sequential Compactness and Total Boundness]
    A metric space $(X, d)$ is sequentially compact if and only if it is complete and totally bounded.
\end{theorem}

\clearpage

\begin{heading}
    Lebesgue number and Lebesgue Covering Lemma
\end{heading}

\vspace{10pt}

\begin{definition}[Lebesgue Number]
    Let $\{G_\alpha : \alpha \in \Lambda\}$ be an open cover of a metric space $(X,d)$. A real number $\lambda > 0$ is said to be a \emph{Lebesgue number} for the open cover $\{G_\alpha\}$ if for each subset $A$ of $X$ with $d(A) < \lambda$, there is at least one set $G_\alpha$ which contains $A$. \\In other words, any subset of $X$ that is "small enough" (in terms of diameter) can be entirely contained within one of the sets in the cover $\mathcal{U}$.
\end{definition}

\clearpage

\begin{theorem}[Lebesgue Covering Lemma]
    Every open cover of a sequentially compact metric space has a Lebesgue number.
\end{theorem}

\clearpage

\begin{heading}
    Sequential Compactness and Compactness
\end{heading}

\vspace{10pt}

\begin{theorem}[In MS, Sequential Compactness $\iff$ Compactness]
    A metric space $(X, d)$ is sequentially compact if and only if it is compact.
\end{theorem}

\clearpage

\begin{theorem}[Continuity $\iff$ Uniform Continuity on Compact Sets]
    Let $(X, d_X)$ and $(Y, d_Y)$ be metric spaces, and let $f : X \to Y$ be a continuous function. If $X$ is compact, then $f$ is uniformly continuous.
\end{theorem}

\clearpage

\begin{heading}
    Equivalent Characterizations of Compactness
\end{heading}

\vspace{10pt}

\begin{theorem}[Equivalent Characterizations of Compactness]
    For a metric space $(X, d)$, the following statements are equivalent:
    \begin{enumerate}
        \item $X$ is compact.
        \item $X$ is sequentially compact.
        \item $X$ is complete and totally bounded.
        \item Bolzano-Weierstrass property
    \end{enumerate}
\end{theorem}

\EndPage
\end{document}